\documentclass[10pt,a4paper]{amsart}
\usepackage[margin=19mm]{geometry}
\usepackage{amsmath,amssymb,mathtools}
\usepackage[scr=rsfs,cal=euler]{mathalfa}
\usepackage{xfrac}
\usepackage[unicode,hidelinks,bookmarks=false]{hyperref}
\newcommand{\ie}{i.e.\ }
\newcommand{\cf}{cf.\ }
\newcommand{\resp}{respectively\ }
\newcommand{\Subsection}{\subsection}
\providecommand{\nobreakdash}{\nobreak\hbox{-}}
\setlength{\emergencystretch}{2em}
\multlinegap0pt
\usepackage{mdframed}
\usepackage{mdframed}
\usepackage{mdframed}
\usepackage{mdframed}
\usepackage[shortlabels]{enumitem}
\usepackage{mathtools}
\usepackage{amssymb}
\usepackage{xcolor}
\usepackage{tikz}
\usepackage{mdframed}
\newtheorem{example}{Example}
\newtheorem{remark}{Remark}
\newtheorem{lemma}{Lemma}
\newtheorem{proposition}{Proposition}
\def\AA{\mathcal{A}}
\def\C{\mathbb{C}}
\def\d{\partial}
\def\g{\mathfrak{g}}
\def\h{\mathfrak{h}}
\def\l{\mathfrak{l}}
\def\o{\mathfrak{o}}
\def\ol{\overline}
\def\p{\mathfrak{p}}
\def\q{\mathfrak{q}}
\def\s{\mathfrak{s}}
\title{A Queer Kac--Moody Construction}
\author{Alexander Sherman, Lior Silberberg}
\date{Source: arXiv:2309.09559v3 (CC BY 4.0)}
\begin{document}
\maketitle

\noindent\emph{Source and licence.} A Queer Kac--Moody Construction. Version arXiv:2309.09559v3 (CC BY 4.0), distributed by arXiv under the Creative Commons Attribution 4.0 International licence, http://creativecommons.org/licenses/by/4.0/, as recorded in the arXiv licence metadata for this exact version and shown on the abstract page https://arxiv.org/abs/2309.09559v3. Full licence notice: \texttt{licenses/arxiv-2309.09559-CC-BY-4.0.txt}.\par\smallskip

\noindent\textit{Editorial scope.} Counted unit: the article's proposition propistion-two-sq2 (source0.tex lines 1579--1591). The article's complete original proof of it is delivered with it (source0.tex lines 1593--1626, one \texttt{proof} environment of 34 lines). No mathematics is written, repaired or re-derived: the statement and the proof are the article's own lines, in the article's own notation, and no retained line is substituted or re-flowed.  Retained uncounted as the source-local context the proof needs: lines 852--861; lines 1278--1335; lines 1395--1408; lines 1559--1568.  Nothing was omitted from any retained span, so the package carries no omission record.  No retained line contains a citation, so the wrapper carries no bibliography and no bibliography entry.  No retained line contains a figure environment, an image-inclusion command or drawing code, so no image is delivered and the package declares no asset dependency.  Editorial formatting only, in addition: an AMS \texttt{amsart} preamble that restates the article's own declarations and macros used by the retained lines, namely the environments \texttt{example}, \texttt{remark}, \texttt{lemma}, \texttt{proposition} on separate counters, together with the standard packages those lines need.  The retained environments print in this document as Example 1, Remark 1, Lemma 1, Proposition 1 in order of appearance, whereas the article prints the same items with the numbering of its own full document.  The complete native source of the article as distributed in the arXiv e-print for this exact version is bundled unchanged as \texttt{sources/147169/source0.tex}.
	\begin{example}\label{example_ts_pts}
		For examples of superalgebras that are not Clifford Kac--Moody, let $\s$ be a symmetrizable Kac--Moody Lie superalgebra, and write $\mathfrak{t}\s$ for the subalgebra of $T\s$ spanned by $\s\otimes\C[\xi]$ and $c$.  It $\ol{\mathfrak{t}}$ is a maximal torus of $\s$, then we set $\h=\ol{\mathfrak{t}}\otimes\C[\xi]\oplus\C\langle c\rangle$ to obtain a self-normalizing quasi-toral subalgebra.  However in this case the root spaces of simple roots, spanned by $e_i\otimes 1$ and $e_i\otimes\xi$ for a Chevalley generator $e_i$ of $\s$, will not be irreducible $\h$-modules because we no longer have the derivation $\d_{\xi}$.  Thus $\mathfrak{t}\s$ is not Clifford Kac--Moody (however it does arise naturally from our more general construction, as $\mathfrak{t}\s$ contains no nontrivial ideals that intersect $\h$ trivially).  
		
		Another example is obtained by considering $\p\mathfrak{t}\s:=\mathfrak{t}\s/(c)\cong\s\otimes\C[\xi]$; and indeed if $\mathfrak{s}$ is not symmetrizable then this is the natural superalgebra to consider.  This case is again clearly not Clifford Kac--Moody, but once again arises naturally from our construction.
		
		Finally, we note that 
		\[
		\mathfrak{t}\s\l(2)\cong\s\q(2) \ \text{ and } \ \mathfrak{pt}\s\l(2)\cong\p\s\q(2).
		\]
	\end{example}

	\begin{remark}
		It follows from our work that if $\AA$ is QKM, then for a simple root $\alpha$, the subalgebra $\g\langle\alpha\rangle$ is isomorphic to a central quotient of $\mathfrak{t}\s$ for a Kac--Moody Lie superalgebra $\mathfrak{s}$ with one simple root (see Example \ref{example_ts_pts}).  Thus QKM algebras are built by `glueing together' extensions of Takiffs of Kac--Moody superalgebras with one simple root.
		
		A trivial way to glue together such Takiffs is via the Takiff construction applied to an arbitrary Kac--Moody Lie superalgebra; however, in a sense, this is a less interesting type of glueing.  We will find that if we avoid such constructions, we are led to more `interesting' Lie superalgebras such as $\q(n)$, and also that the theory becomes very rigid.
	\end{remark}
	
	\subsection{Cartan datum for QKM algebras}\label{section-notation-QKM}
	
	Because QKM algebras only have simple roots of type $Tak(\s)$ for $\s=\s\l(2),\o\s\p(1|2)$, or $\s\l(1|1)$, we can more explicitly present their Cartan datum.  We do that now, while also introducing notation that will be used henceforth throughout the paper.
	
	Let $\AA$ be a QKM Cartan datum, and let $\Pi=\{\alpha_1,\dots,\alpha_n\}$ be the simple roots.  For each $i$, we have the subalgebra $\g\langle\alpha_i\rangle$ generated by $\g_{\alpha_i}$ and $\g_{-\alpha_i}$.  We choose homogeneous bases of $\g_{\pm\alpha_i}$, using the identification $\g_{-\alpha_i}=\g_{\alpha_i}^\vee$, as follows:
	\begin{enumerate}
		\item if $\alpha_i$ is of type $Tak(\s\l(2))$, set $e_i$ and $E_i$ to be nonzero even and odd elements of $\g_{\alpha_i}$, and let $f_i:=e_i^\vee$ and $F_i:=\sqrt{-1}E_i^\vee$.  Rescale the choice of $e_i$ so that $\alpha_i([e_i,f_i])=2$;
		\item if $\alpha_i$ is of type $Tak(\o\s\p(1|2))$ or $Tak(\s\l(1|1))$, set $E_i$ and $e_i$ to be nonzero even and odd elements of $\g_{\alpha_i}$ (notice the change in order), and let $f_i:=\sqrt{-1}e_i^\vee$ and $F_i:=E_i^\vee$.  If $\alpha_i$ is of type $Tak(\o\s\p(1|2))$, then rescale the choice of $e_i$ so that $\alpha_i([e_i,f_i])=1$.
	\end{enumerate}
	It follows from the above choices that:
	\begin{enumerate}
		\item if $\alpha_i$ is of type $Tak(\s\l(2))$, then $e_i,[e_i,f_i],f_i$ form an $\s\l(2)$ triple;
		\item if $\alpha_i$ is of type $Tak(\o\s\p(1|2))$, then $[e_i,e_i],e_i,[e_i,f_i],f_i,[f_i,f_i]$ form an $\o\s\p(1|2)$-quintuple;
		\item if $\alpha_i$ is of type $Tak(\s\l(1|1))$, then $e_i,[e_i,f_i],f_i$ form an $\s\l(1|1)$-triple.
	\end{enumerate}
	On should view the $e_i,E_i,f_i,$ and $F_i$ as the `Chevalley generators' for a QKM algebra.
	
	\subsubsection{Uniqueness of generators}\label{section_uniqueness_gens} Before going further, we note that the above choices are not unique: in particular we may simultaneously rescale $e_i$ and $f_i$ by $(-1)$, or simultaneously rescale $E_i$ and $F_i$ by $\lambda\in\C^\times$. For $Tak(\s\l(1|1))$ we may also rescale $e_i,f_i$ by any $\lambda\in\C^\times$.
	
	\subsubsection{Pure coroots} Set 
	\[
	H_i:=[E_i,f_i];
	\]
	one may check that $H_i\neq0$, and we also always have 
	\[
	H_i=[e_i,F_i].
	\]
	Further set:
	\[
	h_i:=[e_i,f_i] \ \ \ \ c_i:=[E_i,F_i].
	\]
	
	The following table summarizes some of the important properties, which are easy to prove, of the pure coroots we have introduced.
	\begin{center}
		\begin{tabular}{|c|c|c|c|}
			\hline
			& $Tak(\s\l(2))$ & $Tak(\o\s\p(1|2))$ & $Tak(\s\l(1|1))$\\
			\hline
			$H_i^2=$ & $c_i$ & $c_i$ & 0\\
			\hline 
			$c_i$ lies in... & $\s\l(1|1)$-triple & even Heisenberg triple & even Heisenberg triple\\
			& $\alpha_i(c_i)=0$ & $\beta(c_i)=0$ for all roots $\beta$ & $\beta(c_i)=0$ for all roots $\beta$\\
			\hline
			$h_i$ lies in... & $\s\l(2)$-triple & $\o\s\p(1|2)$ quintuple & $\s\l(1|1)$-triple\\
			& $\alpha_i(h_i)=2$ & $\alpha_i(h_i)=1$ & $\alpha_i(h_i)=0$\\
			\hline
		\end{tabular}	
	\end{center}
	
	\begin{remark}\label{remark_root_subalg_really_takiff}
		A more concrete viewpoint on the elements described above is given using the fact that each $\g\langle\alpha_i\rangle$ is a central quotient of $\mathfrak{t}\s=\s\otimes\C[\xi]\oplus\C\langle c\rangle$ for $\s=\s\l(2),\o\s\p(1|2)$, or $\s\l(1|1)$.  In these terms, $e_i,f_i$ are the Chevalley generators of $\s$, $h_i=[e_i,f_i]$, $c_i=c$, and $H_i=h_i\otimes \xi$.  
	\end{remark}

	\begin{lemma}\label{lemma_nontrivial_relation_on_ys}
		For two simple roots $\alpha_i,\alpha_j$ of type $Tak(\s\l(2))$ or $Tak(\o\s\p(1|2))$ we have
		\[
		x_{ij}x_{ji}y_{ji}+y_{ij}\alpha_i(h_j)=y_{ji}\alpha_i(h_i),
		\]
		and
		\[
		y_{ij}y_{ji}(\alpha_i(h_i)-\alpha_j(h_j))=y_{ij}^2\alpha_i(h_j)-y_{ji}^2\alpha_j(h_i).
		\]
		In particular if $\alpha_i(h_i)=\alpha_j(h_j)$, then
		\[
		y_{ij}^2\alpha_i(h_j)=y_{ji}^2\alpha_j(h_i).
		\]
	\end{lemma}

	\begin{lemma}\label{lemma-sq2-rep}
		Let \(\mathfrak{g}(\mathcal{A})\) be a QKM algebra and let \(\alpha \in \Pi\) be of type $Tak(\s\l(2))$. Let \(e,E,f,F,h,c,H\) be the spanning set for $\g\langle\alpha\rangle\cong(\p)\mathfrak{sq}(2)$ as defined in Section \ref{section-notation-QKM}. Let \(\beta \in \Pi\) be a simple root different from \(\alpha\). Then:
		\begin{enumerate}
			\item $\g_{\beta-\beta(h)\alpha}\neq0$, and $\g_{\beta+(1-\beta(h))\alpha}=0$;
			\item  if \(\beta(h)=-1\) and \(v \in \mathfrak{g}_{\beta}\), then 
			\[
			[E,v] + [e,[H,v]] =0.
			\]
		\end{enumerate}
	\end{lemma}

	\begin{proposition}\label{propistion-two-sq2}
		Let \(\mathfrak{g}(\mathcal{A})\) be a QKM algebra and let \(\alpha_i, \alpha_j \in \Pi\) be two simple roots type $Tak(\s\l(2))$. If \(\alpha_i\) and \(\alpha_j\) are connected, then exactly one of the following happens:
		\begin{enumerate}[label=(\roman*)]
			\item $\alpha_i$ and $\alpha_j$ are coupled, in which case either:
			\begin{enumerate}
				\item $\alpha_i$ and $\alpha_j$ are $Y$-coupled; or
				\item $\alpha_i$ and $\alpha_j$ are $X$-coupled and $\alpha_i(h_j)\alpha_j(h_i)=4$.
			\end{enumerate}
			\item \(\alpha_i(h_j) = \alpha_j(h_i) = -1\) and \(\alpha_i(c_j)\alpha_j(c_i) \neq 0\);
			\item \(\alpha_i(h_j)=\alpha_j(h_i) = -2\),  and exactly one of \(x_{ij}\), \(x_{ji}\) is zero;
			\item \(\alpha_i(h_j),\alpha_j(h_i) \in \mathbb{Z}_{\leq -2}\) and \(\alpha_i(c_i)\alpha_j(c_i) \neq 0\).
		\end{enumerate}
	\end{proposition}

	\begin{proof}
		Because we assume $\alpha_i,\alpha_j$ are connected, we have $\alpha_i(h_j)\alpha_j(h_i)\neq0$.  
		
		Now because $\alpha_i(h_i)=\alpha_j(h_j)$, Lemma \ref{lemma_nontrivial_relation_on_ys} implies that:
		\begin{align}
			y_{ij}^2\alpha_i(h_j) = y_{ji}^2\alpha_j(h_i).
		\end{align}
		So if \(y_{ij}y_{ji} = 0\), it must be that both \(y_{ij}=y_{ji} = 0\), implying \(\alpha_i\) and \(\alpha_j\) are $Y$-coupled.
		
		If, on the other hand, \(x_{ij}x_{ji} = 0\), then Lemma \ref{lemma_nontrivial_relation_on_ys} gives:
		\begin{align*}
			2y_{ji} = y_{ij}\alpha_i(h_j), \\
			2y_{ij} = y_{ji}\alpha_j(h_i).
		\end{align*}
		Thus \(y_{ij},y_{ji}\) is a non-trivial solution to the linear system defined by the matrix 
		\[
		\begin{bmatrix} 2 & -\alpha_i(h_j)\\ -\alpha_j(h_i) & 2\end{bmatrix},
		\]
		implying that the determinant is 0, which gives
		\[
		\alpha_i(h_j)\alpha_j(h_i) = 4.
		\]
		In particular, this shows that if $\alpha_i$ and $\alpha_j$ are $X$-coupled, then $\alpha_i(h_j)\alpha_j(h_i)=4$, as desired. On the other hand, if \(\alpha_i\) and \(\alpha_j\) are not coupled but we still have $x_{ij}x_{ji}=0$, then exactly one of \(x_{ij}\), \(x_{ji}\) is zero, and we have $\alpha_i(h_j)\alpha_j(h_i)=4$.  It remains to show then that $\alpha_i(h_j)=\alpha_j(h_i)=-2$, or equivalently that $\alpha_i(h_j)\neq-1$ and $\alpha_j(h_i)\neq-1$, which will follow from our final argument.
		
		Thus suppose that \(\alpha_i(h_j) = -1\) and \(\alpha_i\), \(\alpha_j\) are not coupled. Then Lemma \ref{lemma-sq2-rep} implies \([E_j,e_i] + x_{ji}[e_j,E_i] = 0\). As a result
		\begin{align*}
			0 = [f_i,[E_j,e_i] + x_{ji}[e_j,E_i]] = [E_j,-h_i]+ x_{ji}[e_j,-H_i] = (\alpha_j(h_i)+x_{ji}x_{ij})E_j,
		\end{align*}
		and
		\begin{align*}
			0 = [F_i,[E_j,e_i] + x_{ji}[e_j,E_i]] = [E_j,H_i] + x_{ji}[e_j,c_i] = (y_{ij}-x_{ij}x_{ji}y_{ij})e_j.
		\end{align*}
		Because \(\alpha_i\) and \(\alpha_j\) are not coupled, we have \(y_{ij} \neq 0\), so \(x_{ij}x_{ji}=1\) and \(\alpha_j(h_i) = -1\).  This completes our proof.
	\end{proof}

\end{document}
